\documentclass[10pt,a4paper]{amsart}
\usepackage[margin=19mm]{geometry}
\usepackage{amsmath,amssymb,mathtools}
\usepackage[scr=rsfs,cal=euler]{mathalfa}
\usepackage{xfrac}
\usepackage[unicode,hidelinks,bookmarks=false]{hyperref}
\newcommand{\ie}{i.e.\ }
\newcommand{\cf}{cf.\ }
\newcommand{\resp}{respectively\ }
\newcommand{\Subsection}{\subsection}
\providecommand{\nobreakdash}{\nobreak\hbox{-}}
\setlength{\emergencystretch}{2em}
\multlinegap0pt
\usepackage{bm}
\usepackage{bbm}
\usepackage{dsfont}
\usepackage{xspace}
\usepackage{cancel}
\usepackage{mathtools}
\usepackage{amsthm}
\makeatletter
\@ifundefined{thm}{\newtheorem{thm}{Thm}}{}
\@ifundefined{lem}{\newtheorem{lem}[thm]{Lemma}}{}
\makeatother
\def\S{\mathbb{S}}
\def\sgn{{\rm sgn}}
\newcommand{\lr}[1]{\langle #1 \rangle}
\newcommand{\ucr}[1]{\underset{#1}{\circ}}
\newcommand{\up}[1]{{^{#1}{\it \!\Upsilon}}}
\title[Codimension sequence, grade, and generating degree: an opera]{Codimension sequence, grade, and generating degree: an operadic approach}
\author{Y.-H. Bao, D.-X. Fu, J.-N. Xu, Y. Ye, J.J. Zhang, Y.-F. Zhang, Z.-B. Zhao}
\date{Source: arXiv:2504.19168v4 (CC BY 4.0)}
\begin{document}
\maketitle

\noindent\emph{Source and licence.} Codimension sequence, grade, and generating degree: an operadic approach. Version arXiv:2504.19168v4 (CC BY 4.0), distributed under the Creative Commons Attribution 4.0 International licence, as recorded in the arXiv licence metadata for this exact version and shown on the abstract page https://arxiv.org/abs/2504.19168v4. Full rights notice: licenses/arxiv-2504.19168-CC-BY-4.0.txt.\par\smallskip

\noindent\textit{Editorial scope.} Counted unit: the Lem whose source label is \texttt{yylem3.8} in the article (source0.tex lines 1881--1890, source label \texttt{yylem3.8}), delivered together with the article's own complete original proof of it (source0.tex lines 1893--1989, one \texttt{proof} environment of 97 lines). No mathematics is written, repaired or re-derived: statement and proof are the article's own text line for line. Required context retained uncounted: none. Every definition, display and label of those lines is retained unchanged. Omitted from this package and not redistributed: 1--1880, 1891--1892, 1990--4171. Nothing inside a retained excerpt is omitted or altered: every retained body is delivered exactly as the article's lines are written, so no omission receipt of any kind accompanies this package. Citations: the retained excerpts carry no citation command at all, so no bibliography entry of the article is transcribed and no reference is redistributed. Reference adaptation: none was needed and none was made; every reference inside the delivered unit resolves inside it, so no unresolved reference or citation is produced. Printed slips kept literal: no printed word, symbol, hypothesis or number of the counted statement or of its proof was altered. Layout and class: the article's own class is \texttt{amsart} with packages including latexsym, amsxtra, amscd, ifthen, amsfonts, pifont, verbatim, amsmath, graphicx, amsthm, amssymb, url, color, xy; the wrapper is the lane's pinned \texttt{amsart} editorial wrapper at 10pt on A4, which restates the theorem-like environments the retained text needs and adds the article's own macro definitions that the retained text needs (\texttt{\textbackslash S}, \texttt{\textbackslash lr}, \texttt{\textbackslash sgn}, \texttt{\textbackslash ucr}, \texttt{\textbackslash up}), with extra wrapper packages (bm, bbm, dsfont, xspace, cancel, mathtools). The delivered numbering is the unit's own. Assets: no retained span contains a figure environment, a graphics-inclusion command, a PDF-page inclusion, a table or a TikZ picture; no image file is copied into this package, the package declares no asset dependency and no image file is redistributed with it. Licence: version arXiv:2504.19168v4 of this article is distributed under the Creative Commons Attribution 4.0 International licence, as recorded in the arXiv licence metadata for this exact version and shown on the abstract page https://arxiv.org/abs/2504.19168v4. A full rights notice is delivered at licenses/arxiv-2504.19168-CC-BY-4.0.txt. Characters: every retained excerpt is pure ASCII, so the delivered engine, XeLaTeX with its default Latin Modern text font, reports no missing character.
\begin{lem}
\label{yylem3.8}
Let $k\geq 3$ be an odd number. Then 
\begin{align}
\label{E3.8.1}\tag{E3.8.1}
\tau_{2, \cdots, 2}\ast (\sgn(\sigma)\sigma-1_{k+1})
\in \lr{\up{k}(k)}(k+1)
\end{align}
for any $\sigma\in \S_{k+1}$. 
\end{lem}

\begin{proof}
Applying the isomorphism 
$\Psi_{k+1}\colon V_{k+1}\to \Bbbk\S_{k+1}$, we need only 
to show
\[\Psi_{k+1}(\sgn(\sigma)[x_{\sigma^{-1}(1)}, x_{\sigma^{-1}(2)}]
\cdots [x_{\sigma^{-1}(k)}, x_{\sigma^{-1}(k+1)}]
-[x_1, x_2]\cdots[x_{k}, x_{k+1}])\in \lr{\up{k}(k)}(k+1).\]
We first prove the special case when $k=3$. By a direct calculation, 
we have
\begin{align*}
[x_1, x_3]&[x_2, x_4]+[x_1, x_2][x_3, x_4]\\
=& [[x_1, x_3]x_2, x_4]-[x_1, x_3, x_4]x_2
 +[[x_1, x_2]x_3, x_4]-[x_1, x_2, x_4]x_3\\
=& [[x_1x_2, x_3]-x_1[x_2, x_3], x_4]
 -[x_1, x_3, x_4]x_2+[[x_1x_3, x_2]-x_1[x_3, x_2], x_4]
 -[x_1, x_2, x_4]x_3\\
=& [x_1x_2, x_3, x_4]-[x_1, x_3, x_4]x_2
 +[x_1x_3, x_2, x_4]-[x_1, x_2, x_4]x_3.
\end{align*}
It follows that 
\[\tau_{2, 2}\ast (-(23)-1_{4})
=(1_2\ucr{1}\tau_3)\ast((243)+(34))
-(\tau_3\ucr{1}1_2)\ast(1_{4}+(23))\in \lr{\up{3}(3)}(4),\]
where $(23), (243), (34)$ are cycle permutations in 
$\S_{4}$. 

For the general case, we use induction on the inversion 
number $\phi(\sigma)$. By definition, when $\sigma$ is 
written as $(i_1, i_2, \cdots, i_{k+1})$ with 
$i_s=\sigma^{-1}(s)$ for all $s$, $\phi(\sigma)$ is the 
number of pairs $(i_{s},i_{s'})$ where $s<s'$ and $i_s>i_{s'}$. 
(Sometimes this is called the inversion number of 
$\sigma^{-1}$ by some authors).

If $\phi(\sigma)=0$, then \eqref{E3.8.1} holds obviously 
since $\sgn(\sigma)\sigma-(1)=0$.

Assume that \eqref{E3.8.1} holds when $\phi(\sigma)\le l$. 
When $\phi(\sigma)=l+1$, there are two cases. 

Case 1. There exists $s$ with $2s\le k+1$ such that 
$\sigma^{-1}(2s-1)>\sigma^{-1}(2s)$. 
Set \[\sigma'=(\sigma^{-1}(1), \cdots, \sigma^{-1}(2s-2), 
\sigma^{-1}(2s), \sigma^{-1}(2s-1),
\sigma^{-1}(2s+1), \cdots, \sigma^{-1}(k+1)).\] 
Then we have 
$\phi(\sigma')=\phi(\sigma)-1$, and 
\begin{align*}
 \tau_{2, \cdots, 2}\ast (\sgn(\sigma)\sigma-1_{k+1})
=\tau_{2, \cdots, 2}\ast (\sgn(\sigma')\sigma'-1_{k+1})
\in \lr{\up{k}(k)}(k+1)
\end{align*}
by the induction hypothesis.

Case 2. Suppose $\sigma^{-1}(2t-1)<\sigma^{-1}(2t)$
for all $2t\leq k+1$. Then there exists $s$ with 
$2s\le k$ such that $\sigma^{-1}(2s)>\sigma^{-1}(2s+1)$.
Write $\sigma=(i_1, \cdots, i_{k+1})$ and 
$\sigma'=(i_1, \cdots, i_{2s-1}, i_{2s+1}, 
i_{2s}, i_{2s+2}, \cdots, i_{k+1})$. 
Clearly, $\phi(\sigma')=\phi(\sigma)-1$.  
By the proof of the special case $k=3$, we have 
\begin{align*}
\sgn(\sigma)&[x_{i_1}, x_{i_2}]\cdots 
[x_{i_k}, x_{i_{k+1}}]-[x_1, x_2]\cdots[x_{k}, x_{k+1}]\\
=&\sgn(\sigma)[x_{i_1}, x_{i_2}]\cdots 
[x_{i_k}, x_{i_{k+1}}]-\sgn(\sigma')[x_{i_1}, x_{i_2}]
\cdots[x_{i_{2s-1}}, x_{i_{2s+1}}][x_{i_{2s}}, x_{i_{2s+2}}]
\cdots [x_{i_k}, x_{i_{k+1}}]\\
&+\sgn(\sigma')[x_{i_1}, x_{i_2}]
\cdots[x_{i_{2s-1}}, x_{i_{2s+1}}][x_{i_{2s}}, x_{i_{2s+2}}]
\cdots [x_{i_k}, x_{i_{k+1}}]-
[x_1, x_2]\cdots[x_{k}, x_{k+1}]\\
=&\sgn(\sigma)[x_{i_1}, x_{i_2}]\cdots 
([x_{i_{2s-1}}, x_{i_{2s}}][x_{i_{2s+1}}, x_{i_{2s+2}}]
+[x_{i_{2s-1}}, x_{i_{2s+1}}][x_{i_{2s}}, x_{i_{2s+2}}])
\cdots [x_{i_k}, x_{i_{k+1}}]\\
&+\sgn(\sigma')[x_{i_1}, x_{i_2}]
\cdots[x_{i_{2s-1}}, x_{i_{2s+1}}][x_{i_{2s}}, x_{i_{2s+2}}]
\cdots [x_{i_k}, x_{i_{k+1}}]-
[x_1, x_2]\cdots[x_{k}, x_{k+1}]\\
=& \sgn(\sigma)[x_{i_1}, x_{i_2}]
\cdots \left([x_{i_{2s-1}}x_{i_{2s}}, x_{i_{2s+1}}, x_{i_{2s+2}}]
-[x_{i_{2s-1}}, x_{i_{2s+1}},x_{i_{2s+2}}]x_{i_{2s}}\right.\\
&\qquad+[x_{i_{2s-1}}x_{i_{2s+1}}, x_{i_{2s}}, x_{i_{2s+2}}]
-[x_{i_{2s-1}}, x_{i_{2s}},x_{i_{2s+2}}]x_{i_{2s+1}}\left. \right)
\cdots[x_{i_k}, x_{i_{k+1}}]\\
&+\sgn(\sigma')[x_{i_1}, x_{i_2}]\cdots[x_{i_{2s-1}}, x_{i_{2s+1}}]
[x_{i_{2s}}, x_{i_{2s+2}}]\cdots [x_{i_k}, x_{i_{k+1}}]-
[x_1, x_2]\cdots[x_{k}, x_{k+1}].
\end{align*}
Therefore, from the induction hypothesis, it follows that 
\[\Psi_{k+1}(\sgn(\sigma)[x_{\sigma^{-1}(1)}, x_{\sigma^{-1}(2)}]
\cdots [x_{\sigma^{-1}(k)}, x_{\sigma^{-1}(k+1)}]-[x_1, x_2]
\cdots[x_{k}, x_{k+1}])\in \lr{\up{k}(k)}(k+1).\]
This completes the proof.
\end{proof}

\end{document}
